\documentclass[11pt,a4paper]{article}
\usepackage[margin=25mm,headheight=14pt]{geometry}
\usepackage[T1]{fontenc}
\usepackage{lmodern,amsmath,amssymb,amsthm,mathtools}
\usepackage{booktabs,array,enumitem,microtype,needspace}
\usepackage[dvipsnames]{xcolor}
\usepackage{hyperref,fancyhdr}
\hypersetup{colorlinks=true,linkcolor=MidnightBlue,citecolor=MidnightBlue,
urlcolor=MidnightBlue,pdftitle={Modularity of elliptic curves over CM and totally real fields},
breaklinks=true}
\Urlmuskip=0mu plus 1mu
\setlength{\parindent}{0pt}
\setlength{\parskip}{4.5pt}
\setlist{topsep=4pt,itemsep=2pt,parsep=0pt}
\pagestyle{fancy}\fancyhf{}
\fancyhead[L]{\small Elliptic curves over CM and totally real fields}
\fancyhead[R]{\small Research draft}
\fancyfoot[C]{\thepage}
\newtheorem{theorem}{Theorem}[section]
\newtheorem{proposition}[theorem]{Proposition}
\newtheorem{lemma}[theorem]{Lemma}
\newtheorem{corollary}[theorem]{Corollary}
\theoremstyle{definition}\newtheorem{definition}[theorem]{Definition}
\theoremstyle{remark}\newtheorem*{remark}{Remark}
\newcommand{\Q}{\mathbb Q}\newcommand{\Z}{\mathbb Z}
\newcommand{\C}{\mathbb C}\newcommand{\A}{\mathbb A}
\newcommand{\F}{\mathbb F}\newcommand{\OO}{\mathcal O}
\newcommand{\GL}{\operatorname{GL}}\newcommand{\SL}{\operatorname{SL}}
\newcommand{\PSL}{\operatorname{PSL}}\newcommand{\PGL}{\operatorname{PGL}}
\newcommand{\Gal}{\operatorname{Gal}}\newcommand{\Res}{\operatorname{Res}}
\newcommand{\Jac}{\operatorname{Jac}}\newcommand{\Spec}{\operatorname{Spec}}
\newcommand{\diag}{\operatorname{diag}}\newcommand{\tr}{\operatorname{tr}}
\newcommand{\WD}{\operatorname{WD}}\newcommand{\BC}{\operatorname{BC}}
\newcommand{\ad}{\operatorname{ad}}\newcommand{\Inf}{\operatorname{Inf}}
\newcommand{\im}{\operatorname{im}}\newcommand{\ur}{\mathrm{ur}}
\newcommand{\ord}{\mathrm{ord}}\newcommand{\e}{\epsilon}
\newcommand{\HH}{\mathrm H^1_{\mathrm{\acute et}}}
\newcommand{\srcM}{\cite{M}}
\newcommand{\Ginput}{\mathcal G}
\newcommand{\Tinput}{\mathcal T_5}
\title{Modularity of elliptic curves\\over CM and totally real fields}
\author{AI-assisted research draft\thanks{AI assistance: OpenAI ChatGPT was used for mathematical exploration and drafting; OpenAI Codex was used for editing and compilation.}}
\date{8 October 2026}
\hypersetup{pdfsubject={Research draft}}
\begin{document}
\maketitle

\begin{abstract}
We assemble the field extension of the geometric prime switch in
\texttt{openai/math}, family 030, and a proposed treatment of its remaining
characteristic-five obstruction. The source's explicit geometric and local
inputs imply modularity over every CM field not containing a primitive fifth
root of unity, with no Galois or degree restriction. Quadratic descent then
gives modularity over every totally real field. For CM fields containing fifth
roots of unity, a split Tate condition with parameter valuation prime to five
excludes the exceptional adjoint cohomology class. We give a partially framed
ordinary patching argument using the local rings of Allen--Khare--Thorne.
The width-three second degeneration of the geometric family produces precisely
this Tate condition. Validation of the new patching argument would complete
the extension to all CM fields. Independently of the geometric construction,
it would remove the fifth-root-of-unity restriction from Caraiani--Newton's
residual-mod-three modularity criterion.
\end{abstract}

\noindent\fbox{\parbox{\dimexpr\linewidth-2\fboxsep-2\fboxrule\relax}{
\textbf{Status.} This is a research draft. The package $\Ginput$ in
Section~\ref{sec:package} records the geometric and local assertions used from
family 030; that source manuscript is not independently certified here.
The argument $\Tinput$ in Proposition~\ref{prop:global} is a new patching
adaptation presented for independent review. The CM theorem without fifth
roots of unity and the totally real theorem use $\Ginput$ only. The all-CM
extension uses both $\Ginput$ and $\Tinput$.
}}

\section{Main statements and conventions}
\label{sec:statements}

A CM field is a totally imaginary quadratic extension of a totally real number
field. It need not be Galois over $\Q$. For an elliptic curve $E/K$ we write
\[
 V_\ell(E)=\HH(E_{\overline K},\Q_\ell),\qquad
 \det V_\ell(E)=\e_\ell^{-1},
\]
with labelled Hodge--Tate weights $\{0,1\}$. We use weight zero for the
cohomological automorphic normalization; over a totally real field this means
parallel classical weight two. The notation $E[\ell]$ uses the usual Tate-module
convention. Passing to its dual does not change any residual irreducibility,
image-up-to-conjugacy, or decomposed-generic condition used below.

For a non-CM elliptic curve, modularity means a cuspidal weight-zero automorphic
representation with its Galois representations and finite local parameters.
For a CM elliptic curve we use the classical Hecke-character construction: its
realization is an isobaric sum when the CM field is contained in the ground
field, and quadratic automorphic induction otherwise. Thus cuspidality is not
asserted for a split CM realization.

For a residual representation of characteristic $\ell$, decomposed genericity
means that there is a rational prime $q\ne\ell$, completely split in the
ground field, such that the representation is unramified at every place above
$q$ and every ratio of two distinct indexed Frobenius eigenvalues differs from
$q$ modulo $\ell$; see \cite{CN}, Definition 2.1.27. We will always construct
witnesses with $q\equiv1\pmod\ell$ and two distinct eigenvalues.

Here are the conclusions and their dependencies. The proofs appear in the
indicated sections; the input symbols are defined there without assuming any
all-CM or all-totally-real modularity statement.
\begin{center}
\begin{tabular}{@{}p{.47\linewidth}p{.17\linewidth}p{.22\linewidth}@{}}
\toprule
Conclusion & Inputs & Location\\
\midrule
All elliptic curves over CM $K$ with $\zeta_5\notin K$
 & $\Ginput$ & Theorem~\ref{thm:CM}\\[3pt]
All elliptic curves over totally real fields
 & $\Ginput$ & Corollary~\ref{cor:TR}\\[3pt]
All elliptic curves over every CM field
 & $\Ginput+\Tinput$ & Corollary~\ref{cor:allcm}\\[3pt]
Residual-mod-three modularity over every CM field
 & $\Tinput$ & Corollary~\ref{cor:mod3}\\
\bottomrule
\end{tabular}
\end{center}
The last row retains decomposed genericity and absolute irreducibility over
the field obtained by adjoining $\zeta_3$; it is not an unrestricted modularity
claim. In particular the totally real conclusion is logically independent of
the proposed exceptional-five step.

The two field choices driving the argument are different. First prove the
auxiliary elliptic curve modular over its parameter's ground field; only then
adjoin the cyclotomic characters for the three congruences. For a totally real
ground field, pass to a CM quadratic extension chosen to avoid $\zeta_5$ and
descend. For the remaining CM case, keep the characteristic-five seed and
control its local ordinary deformation condition geometrically.

\section{Geometric systems and local inputs}
\label{sec:package}

We write \(G_L=\Gal(\overline L/L)\) and use the cohomological convention
\(\det\HH(E,\Q_\ell)=\epsilon_\ell^{-1}\), with labeled Hodge--Tate weights
\(\{0,1\}\). The following package, denoted \(\Ginput\), records the source inputs.
The source [M] is the manuscript dated 4 October 2026 listed in the bibliography;
its assertions are used in precisely the forms stated below.

\begin{enumerate}[label=\textbf{G\arabic*.},leftmargin=40pt]
\item \textbf{Covers and compatible systems.}
For an elliptic curve $E/k$ over a number field and odd $n$, \srcM,
Lemma 3.1 constructs $U\subset\A^1_k$, quadratic twists $B_t$ of $E$, and
smooth genus-$m$ curves $C_m(t)$ for $m\mid n$. The covers have rotations
after adjoining $\mu_m$, and an involution $w$ over $k$ inverting those
rotations. For $n=3p$, the character pieces $V_\theta$ over
$F=k(\mu_{3p})$ are rank-two compatible systems. Proposition 5.1 gives
semisimplicity, determinant $\e_\ell^{-1}$, labelled Hodge--Tate weights
$\{0,1\}$ at every embedding, potential semistability, and compatible
Frobenius polynomials. The construction is over arbitrary number fields;
its Hodge-number calculation is performed embedding by embedding.

\item \textbf{Integral character congruences.}
Proposition 4.2 gives the rotation-module isomorphism
\[
\mathrm H^1(C_m,\Z_\ell)\simeq\Z_\ell[\mu_m]^2,
\]
including when $\ell\mid m$. Thus evaluation at characters with the same
reduction gives isomorphic semisimplified residual representations, as in
Lemma 5.3 and Corollary 5.4. In particular the congruences used below do not
require integral character projectors with denominator divisible by $3p$.

\item \textbf{Two geometric monodromy elements.}
For a nontrivial residual character with value $\lambda$, Theorem 4.5 gives,
in one basis, both matrices
\[
T_a=\begin{pmatrix}1&\sigma c_\lambda\\0&1\end{pmatrix},\qquad
T_b=\begin{pmatrix}1&0\\\tau c_\lambda&1\end{pmatrix},\qquad
c_\lambda=2-\lambda-\lambda^{-1},\quad\sigma,\tau\in\{1,-1\}.
\]
The two loops act simultaneously on all character pieces of each cover.

\item \textbf{Compact-type local specialization.}
The first degeneration has Jacobian $E^n$. Proposition 7.5 and Corollary 7.6
provide nonempty local opens at every $v\mid3p$, over arbitrary finite
extensions of the relevant $\Q_\ell$, in which every rank-two piece has the
same type as $E$: potentially crystalline ordinary, potentially crystalline
with slopes $1/2$, or potentially totally toric with nonzero monodromy.

\item \textbf{The second boundary model.}
On the constant-$E$ parameter cover, let $R_*$ be a boundary point with
$Q=[n]R_*\in E[2]\setminus\{O\}$ and $O,\pm R_*,Q$ distinct. The local
blowup and normalization calculation of \srcM, Lemma 3.2, equation (3.7),
and Lemma 3.5 gives an unramified degree-three cover of $E$, possibly
disconnected, with Kummer class $-2(n/3)R_*$ and a geometrically rational tail,
meeting in three geometric nodes of smoothing order one. The involution
extends to this model. Section~\ref{sec:tategeometry} proves connectedness
at exact-order-six $R_*$, splitness after adjoining $E[6]$, and the quotient's
width three; none of those arithmetic conclusions is assumed here.
Only the all-CM extension needs this second boundary input.
\end{enumerate}

\paragraph{Local automorphic comparison (LC).}
The package $\Ginput$ also includes the local-comparison assertion of \srcM,
Lemma 9.2: when a character system has become automorphic, its next local
automorphic input has zero monodromy in the potentially good case and is a
quadratic twist of Steinberg in the toric case, with the ordinarity matching
G4. This assertion uses compatibility at a different coefficient prime, purity,
and genericity of cuspidal $\GL_2$ local components. It is a separate source
input from G1--G5. No degree-two condition is used in the argument stated
there. The simultaneous arithmetic Hilbert specialization over arbitrary CM
fields is proved in Section~\ref{sec:specialization}, rather than assumed
as part of $\Ginput$.

The elliptic identifications from Proposition 5.2 are
\[
 V_1=\HH(B_t),\qquad
 V_\eta=\HH(A_t),\qquad
 A_t=\Jac(C_3(t)/\langle w\rangle),\qquad \operatorname{ord}(\eta)=3.
\]
The quotient and its Jacobian are defined over \(k\), because \(C_3(t)\) and
\(w\) are. The fact that the source introduces \(A_t\) over \(F\) does not
prevent this earlier descent. A rational point on the genus-one quotient is
unnecessary for its Jacobian to be an elliptic curve.

We also use Serre's open-image theorem \cite{Serre}, Hilbert irreducibility
with weak approximation, Chebotarev, and automorphic Galois representations.
The exact base-change and descent result is \cite{ACC}, Proposition 6.5.13,
which permits both imaginary CM and totally real fields and does not require
a polarization hypothesis. Its proof uses cyclic base change and descent
from \cite{AC}, Chapter 3, Theorems 4.2 and 5.1. The specific automorphy inputs
are Caraiani--Newton \cite{CN}, Theorem 5.2, Theorem 6.1, and Lemma 6.1.3.
Theorem 5.2 allows arbitrary CM fields and ramification at its coefficient
prime. Its extra condition in characteristic five will not occur in the three
lifts.

\section{Normal closures and large residual images}

\begin{lemma}\label{lem:cm}
The normal closure of a CM field is CM, with central complex conjugation.
Every subfield of a CM field is either totally real or CM. Consequently, for
a CM field \(K\), the extension \(K(\mu_N)/K\) is abelian and all its
intermediate fields are CM.
\end{lemma}
\begin{proof}
If \(c_K\) is the CM involution, then every embedding \(\sigma:K\to\C\)
satisfies \(\overline{\sigma(x)}=\sigma(c_Kx)\).
The normal closure is generated by the conjugates of \(K\), so complex
conjugation commutes with every automorphism there. Its fixed field is totally
real: every conjugate of an element fixed by this central involution is real.
The involution is nontrivial, so the normal closure is CM. The same argument
applies to composita of CM fields.

Inside a Galois CM normal closure, central complex conjugation preserves every
subfield. Its restriction either acts trivially, giving a totally real field,
or makes the subfield a quadratic totally imaginary extension of its totally
real fixed field. Apply this to the subfields of the compositum
\(K(\mu_N)\). A field containing \(K\) cannot be totally real. Finally,
\(K(\mu_N)/K\) is abelian because it is obtained from the abelian extension
\(\Q(\mu_N)/\Q\) by base extension.
\end{proof}

Let now \(K\) be any CM field, \(d=[K:\Q]\), and let \(\widetilde K\) be its
\emph{fixed} normal closure. For a non-CM elliptic curve \(E/K\), Serre's theorem
allows a prime
\[
 p>\max(100,2d+1),\qquad
 \rho_{E,p}(G_{\widetilde K})\supseteq\SL_2(\F_p).
\]
Set
\[
 F=K(\mu_{3p}),\qquad D_0=\widetilde K(\mu_{3p}).
\]
The automorphy field will be \(F\), not the normal closure. The latter is used
only to control Frobenius conditions at every conjugate.

\begin{lemma}\label{lem:survives}
For every embedding \(\sigma:K\hookrightarrow\widetilde K\),
\(\rho_{E^\sigma,p}(G_{D_0})=\SL_2(\F_p)\).
\end{lemma}
\begin{proof}
Normality of \(\widetilde K/\Q\) transports the large-image assertion to every
conjugate \(E^\sigma\). Since \(D_0/\widetilde K\) is abelian, the image of
\(G_{D_0}\) contains the commutator subgroup of the image of
\(G_{\widetilde K}\). The group \(\SL_2(\F_p)\) is perfect for \(p>3\);
hence that commutator contains \(\SL_2(\F_p)\). Conversely, the determinant is
trivial over \(D_0\), which contains \(\mu_p\). This proves equality.
\end{proof}
This argument deals with the dependence of \(D_0\) on \(p\). It does not compare
\(p\) with the degree of a field chosen after \(p\), and does not require a
large image after passing to the finite extension cut out by the later
Hilbert specialization.

\begin{lemma}\label{lem:genericE}
For every quadratic twist \(B/K\) of \(E\), its residual representation at \(p\)
is absolutely irreducible and decomposed generic over \(F\).
\end{lemma}
\begin{proof}
On \(G_{D_0}\), form the joint image
\[
 J\subseteq\prod_{\sigma:K\hookrightarrow\widetilde K}\SL_2(\F_p).
\]
Each of its \(d\) projections is surjective. In \(\SL_2(\F_p)\), precisely
\(p^2\) elements have trace 2: writing such a matrix as
\(\left(\begin{smallmatrix}a&b\\c&2-a\end{smallmatrix}\right)\), its determinant
condition is \(bc=-(a-1)^2\). For \(a=1\) there are \(2p-1\) solutions;
the other \(p-1\) values give \(p-1\) solutions each. The count for trace
\(-2\) is the same.

The fibers of each surjective projection from \(J\) have equal cardinality.
The union bound therefore gives
\begin{equation}
 \frac{\#\{g\in J:\text{some factor has trace }2\text{ or }-2\}}{|J|}
 \leq \frac{2dp}{p^2-1}<1.                     \label{eq:union-bound}
\end{equation}
No independence of the conjugate residual representations is assumed.
Choose an element regular semisimple in every factor. In a finite Galois
extension of \(\Q\) recording \(D_0\) and the torsion fields, its conjugacy
class still has this property, because conjugation permutes the \(d\) factors.

Chebotarev gives rational primes \(q\) completely split in \(D_0\), hence in
\(F\), with these Frobenius actions. In particular \(q\equiv1\pmod p\).
At every place over \(q\), distinct eigenvalues give the decomposed-generic
condition: neither eigenvalue ratio equals \(q\) modulo \(p\).
Absolute irreducibility over \(F\) follows from the full special-linear image
on \(G_{D_0}\subseteq G_F\). Quadratic twisting preserves irreducibility and
eigenvalue ratios; exclude the finitely many primes ramified in the twist.
\end{proof}

\section{Simultaneous arithmetic specialization}\label{sec:specialization}

\begin{lemma}\label{lem:specialization}
Assume \(\Ginput\). Prescribe nonempty opens \(\Omega_v\subset U(K_v)\) at
finitely many finite places, including every \(v\mid3p\). There is
\(t\in U(K)\) in all these opens for which, over \(F\):
\begin{enumerate}
\item every nontrivial character space for
\((m,\ell)=(3,p),(3,5),(p,3)\) is absolutely irreducible on
\(G_{F(\zeta_\ell)}\) and decomposed generic;
\item the same assertions hold for \(V_1\bmod p\);
\item \(A_t\) is non-CM, and the image of \(A_t[5]\) on \(G_D\)
is \(\SL_2(\F_5)\), for the auxiliary Galois field \(D\) used below.
In particular, if \(\mu_5\subset K\), its full mod-five image is
\(\SL_2(\F_5)\).
\end{enumerate}
\end{lemma}
\begin{proof}
Put \(W=\Res_{K/\Q}U\), an open of \(\A^d_\Q\), and choose a finite Galois
field \(D/\Q\) containing \(D_0\) and \(\zeta_5\); \(D=D_0(\zeta_5)\) suffices.
Then
\[
 W_D\simeq\prod_{\sigma:K\hookrightarrow D} U^\sigma_D.
\]
For each listed pair, pull back the finite local system
\(\mathrm H^1(C_m,\F_\ell)\) from every factor. Intersect the kernels of these
finite actions in \(\pi_1(W_D)\), and take the normal core in \(\pi_1(W)\).
We obtain a finite quotient
\[
 \pi_1(W)\twoheadrightarrow H\twoheadrightarrow\Gal(D/\Q).
\]

In one geometric factor, the two opposite nonidentity unipotents in G3 have
different fixed lines, hence force absolute irreducibility. There is also a
single geometric element regular semisimple on \emph{all} its nontrivial
characters for each pair. For \(m=3\), \(c_\lambda=3\), and \(T_aT_b^2\)
has trace \(20\) or \(-16\), neither equal to \(\pm2\) modulo 5 or modulo
\(p>100\). For \(m=p,\ell=3\), use \(T_aT_b\), with trace
\(2+\sigma\tau c_\lambda^2\). A failure of regular semisimplicity would imply
\(c_\lambda^2=\pm4\), giving a polynomial of degree at most four for
\(\lambda\) over \(\F_3\). Its order \(p\) would divide \(3^e-1\) for some
\(e\leq4\), impossible for \(p>100\).

Geometric loops in different factors act independently on the systems pulled
back from those factors. Multiplying the chosen factor loops gives, for each
pair, an element \(h_{m,\ell}\) of the constant-field kernel which is regular
semisimple on every factor and character. The opposite unipotents are also
in this kernel.

For each rational prime \(q\),
\[
 W(\Q_q)=\prod_{v\mid q}U(K_v).
\]
Thus the local conditions are nonempty opens on \(W(\Q_q)\). Hilbert
irreducibility with weak approximation produces \(x\in W(\Q)\) in those
opens with surjective specialization \(s_x:G_\Q\to H\). This is precisely
the arbitrary-dimensional Hilbert statement already recorded as \srcM,
Lemma 8.3. The constant-field map implies
\begin{equation}
 s_x(G_D)=\ker(H\to\Gal(D/\Q)).                 \label{eq:specialization-kernel}
\end{equation}
The point \(x\) corresponds to one \(t\in U(K)\); its coordinates over \(D\)
are the actual conjugates \(\sigma(t)\), not independently selected parameters.

Equation~\eqref{eq:specialization-kernel} retains the opposite unipotents over \(G_D\), hence also over
\(G_{F(\zeta_\ell)}\). Since \(\ell\nmid m\) for the three pairs, character
decomposition of the smaller cover commutes with reduction. The comparison
with the \(C_{3p}\)-pieces preserves semisimplified reductions.

Conjugation permutes the specialized conjugate covers and their rotation
characters. Thus every conjugate of \(h_{m,\ell}\) has the required
regular-semisimple action. Chebotarev in the extension cut out by \(s_x\)
gives rational primes split completely in \(D\) and regular semisimple at
every place above them. They are \(1\bmod\ell\), establishing decomposed
genericity. Each pair may use a different prime. Assertion 2 is
Lemma~\ref{lem:genericE}.

For $m=3$ and $\ell=5$, $c_\lambda=3$. Powers of the two nonzero
unipotents generate both root subgroups, hence a group of order $120$
isomorphic to $\SL_2(\F_5)$. The rational isomorphism between the character
system and $V_5(A_t)$ identifies their semisimplified residual representations.
Since the character representation is absolutely irreducible, the actual
residual representation on $A_t[5]$ is isomorphic to it. In an $\F_5$-basis
of $A_t[5]$ this subgroup has determinant one and order $120$, so it is
$\SL_2(\F_5)$. The determinant is trivial on $G_D$; if $\mu_5\subset K$
it is trivial on $G_K$ as well. This proves the asserted full-image statements
without requiring an integral isomorphism of the two original lattices.

Finally, \(A_t\bmod p\) has the two opposite noncommuting unipotents of order
\(p\). A CM elliptic curve has all endomorphisms defined over an extension
of degree at most two; its residual image consequently has an abelian normal
subgroup of index at most two. Both odd-order unipotents would lie in that
subgroup and commute, a contradiction.
\end{proof}

\section{The seed field, three lifts, and descent}

\begin{theorem}[Conditional CM scope extension]\label{thm:CM}
Assume \(\Ginput\). Let \(K\) be any CM number field with
\(\zeta_5\notin K\). Every elliptic curve \(E/K\) is modular.
For non-CM \(E\), there is a cuspidal weight-zero representation
\(\pi_E\) of \(\PGL_2(\A_K)\) realizing its cohomology at every coefficient
prime, with full finite-place local compatibility in the usual
Tate normalization.
\end{theorem}
\begin{proof}
First let \(E\) be non-CM. Choose \(p,F,t\) as above, using all the opens
from G4. The auxiliary elliptic curve \(A_t\) is over \(K\).
Its residual representation modulo 5 is absolutely irreducible on
\(G_{F(\zeta_5)}\), hence on \(G_{K(\zeta_5)}\).
The rational split prime witnessing decomposed genericity over \(F\) also
witnesses it over \(K\): all relevant local residue degrees are one.

Caraiani--Newton, Theorem 6.1, applies to \(A_t/K\).
\emph{This is the only step requiring \(\zeta_5\notin K\).}
Base change its modular representation through the abelian extension \(F/K\),
using \cite{ACC}, Proposition 6.5.13(1).
It remains cuspidal. Indeed, in a cyclic prime-degree step, noncuspidal base
change would give a nontrivial self-twist by a character of that cyclic
extension. A corresponding intertwiner of Galois representations would be
scalar after restriction to the upper field, by absolute irreducibility and
Schur's lemma; a scalar cannot implement a nontrivial twist.
Irreducibility at every step follows already from the residual
irreducibility over \(F\).

Let \(\eta,\psi\) have orders 3 and \(p\). G2 gives
\begin{equation}
 V_\eta\equiv_p V_{\eta\psi}
       \equiv_3 V_\psi
       \equiv_p V_1=\HH(B_t).                  \label{eq:switch}
\end{equation}
Starting with the base change of \(A_t\), apply Caraiani--Newton,
Theorem 5.2, successively as follows:
\begin{center}
\begin{tabular}{lll}
\toprule
Automorphic input&Coefficient prime&Target\\
\midrule
\(V_\eta\)&\(p\)&\(V_{\eta\psi}\)\\
\(V_{\eta\psi}\)&\(3\)&\(V_\psi\)\\
\(V_\psi\)&\(p\)&\(V_1\)\\
\bottomrule
\end{tabular}
\end{center}
Determinants, all labeled weights, and potential semistability are G1.
Residual hypotheses are Lemmas~\ref{lem:genericE} and
\ref{lem:specialization}. In the potentially good cases G4 and LC give matching
ordinarity and zero automorphic monodromy; in the toric case they give
noncrystalline representations and an ordinary Steinberg input.
Compatible Frobenius polynomials transport the automorphic realization to the
coefficient prime needed for the next lift, using semisimplicity and
Chebotarev. Thus all hypotheses in Theorem 5.2 are met.
The lifting primes are \(p,3,p\), so its characteristic-five exception
never enters. In particular, \(\zeta_5\) is allowed to lie in \(F\).

We have modularity of \(B_t/F\). The extension \(F/K\) is abelian, with CM
intermediate fields by Lemma~\ref{lem:cm}. Apply \cite{ACC}, Proposition 6.5.13(2). Its argument descends through
a tower of cyclic prime-degree extensions: For completeness, if \(L'/L\) is one such step and
\(r|_{G_{L'}}\) is realized by a cuspidal representation \(\Pi\), extension of
\(r\) to \(G_L\) gives \(\Gal(L'/L)\)-invariance of \(\Pi\) by strong
multiplicity one. Cyclic descent gives \(\pi/L\). Its complex archimedean
parameters still have weight zero. After identifying restrictions,
Schur's lemma gives
\[
 r_\ell(\pi)=r\otimes\delta,\qquad
 \delta:\Gal(L'/L)\longrightarrow\overline{\Q}_\ell^\times.
\]
Twisting by the inverse Hecke character removes \(\delta\).
Absolute irreducibility holds at every step because it holds over \(F\).
No nonsolvable descent and no prime-to-\(\ell\) extension-degree assumption
is used.

Twist \(B_t/K\) back to \(E/K\). Caraiani--Newton, Lemma 6.1.3, supplies
realization at all coefficient primes and
\begin{equation}
 \WD\bigl(\HH(E_{\overline K},\Q_\ell)|_{G_{K_v}}\bigr)^{F\text{-ss}}
   \simeq \operatorname{rec}^{T}_{K_v}(\pi_{E,v})
   \quad(v\nmid\ell),                          \label{eq:local-compat}
\end{equation}
after the chosen coefficient isomorphism. Choosing \(\ell\) away from each
place yields every finite local factor.

For CM \(E\), the classical Hecke-character construction applies over the
compositum with its imaginary quadratic CM field \(M\). It gives an isobaric
sum if \(M\subseteq K\), and quadratic automorphic induction otherwise.
\end{proof}

\section{Totally real descent}\label{sec:TR}

\begin{corollary}[Conditional totally real modularity]\label{cor:TR}
Assume \(\Ginput\). Every elliptic curve over every totally real number
field \(k\) is modular of parallel weight two, with matching finite local
parameters.
\end{corollary}
\begin{proof}
For non-CM \(E/k\), choose an odd rational prime \(q\ne5\) unramified in \(k\)
and set \(K=k(\sqrt{-q})\). This is a quadratic CM extension, ramified at
places above \(q\). Since \(k(\zeta_5)/k\) is unramified at those places,
\(\zeta_5\) cannot belong to \(K\): otherwise the nontrivial subextension
\(k(\zeta_5)/k\) of this quadratic extension would equal \(K/k\).
Theorem~\ref{thm:CM} gives a modular representation \(\Pi/K\).

Since \(E\) is defined over \(k\), \(\Pi\) is invariant under \(\Gal(K/k)\).
Apply \cite{ACC}, Proposition 6.5.13(2), to \(V_\ell(E)\), choosing \(\ell\)
so that its restriction to \(G_K\) is irreducible. It produces a cuspidal
regular algebraic representation \(\pi/k\) realizing \(V_\ell(E)\).
The real archimedean type requires a check beyond CM-to-CM descent.
At a complex place of \(K\) its unitary parameter is
\[
 \chi\oplus\chi^c,\qquad \chi(z)=z/|z|.
\]
The characters are distinct. Any \(W_\mathbb R\)-parameter restricting to
this must interchange the two lines, and is
\(\operatorname{Ind}_{W_\C}^{W_\mathbb R}\chi\), the parameter of
weight-two discrete series. For example, the nontrivial Weil element can
act by \(\left(\begin{smallmatrix}0&1\\-1&0\end{smallmatrix}\right)\).
There is no principal-series extension of these two angular characters.
Thus $\pi$ has parallel classical weight two. The character correction
implicit in the descent proposition is at most a sign twist at a real place,
which preserves this discrete series. Its central character is trivial:
its base change is trivial, so it is either trivial or the quadratic character
of $K/k$, and the latter has nontrivial sign at every real place whereas
weight-two discrete series has trivial central character.

One can also obtain all finite local parameters using just the CM comparison
already invoked. Fix a finite place \(v\) of \(k\). Choose a fresh finite place
\(w\ne v\), away from 2 and 5, and by weak approximation choose \(a\in k^\times\)
negative at every real place, a square in \(k_v^\times\), and of valuation one
at \(w\). The extension \(L^{(v)}=k(\sqrt a)\) is CM, splits at \(v\), and avoids
\(\zeta_5\) by its ramification at \(w\).
The modular representation of \(E/L^{(v)}\) and \(\BC_{L^{(v)}/k}(\pi)\) agree at
good places, hence everywhere by strong multiplicity one; base change stays
cuspidal by the same irreducibility argument. At \(v\) local base change is
the identity. Formula~\eqref{eq:local-compat} over \(L^{(v)}\) therefore gives the local parameter
over \(k_v\). Repeat for each \(v\). CM curves are covered by classical
quadratic automorphic induction.
\end{proof}


\section{The exceptional characteristic-five obstruction}
\label{sec:exception}

The three lifts in \eqref{eq:switch} use $p,3,p$, with $p>5$. Their lifting
hypotheses allow $\zeta_5$ in the automorphy field. The root-of-unity restriction
in Theorem~\ref{thm:CM} enters solely when the auxiliary curve $A_t$ is first
proved modular by CN, Theorem 6.1. Its mod-five seed proposition, Proposition
6.1.5, has no such restriction; its mod-three seed proposition, Proposition
6.1.6, uses the internal 3--5 switch of AKT, Proposition 9.15.

Suppose $\mu_5\subset K$ and $\bar\rho$ has image $\SL_2(\F_5)$.
For $M=\ad^0\bar\rho$ one has the exceptional line
\begin{equation}\label{eq:lie}
H^1(\SL_2(\F_5),M)\simeq\F_5.
\end{equation}
It is not detected by ordinary Taylor--Wiles primes. The image restrictions in
AKT, Theorem 8.1, and CN, Theorem 5.2, exclude this global obstruction.
Twisting does not change the projective image; soluble restriction preserves
full $\SL_2(\F_5)$ image because that group is perfect. Our alternative is to
exclude the line from the relevant dual Selmer group using a local ordinary
condition. Sections~\ref{sec:local}--\ref{sec:seed} are independent of $\Ginput$.

\section{The local Tate condition}
\label{sec:local}

Let $L/\Q_5$ contain $\mu_5$, and let $T_q/L$ be a Tate elliptic curve with
parameter $q$. Assume
\begin{equation}\label{eq:tate}
 v_L(q)>0,\qquad 5\nmid v_L(q).
\end{equation}
Up to a choice of basis, and using either the Tate-module convention or its
cohomological twist, its residual representation has the form
\[
 \bar\rho_q(g)=\begin{pmatrix}1&a_q(g)\\0&1\end{pmatrix}.
\]
The character $a_q$ is the Kummer character of $q$, after choosing an
identification $\mu_5\simeq\F_5$. Since $5\nmid v_L(q)$,
$L(q^{1/5})/L$ is totally ramified of degree 5. In particular the residual
inertia image is the full unipotent subgroup $U\simeq C_5$.

Write
\[
 e=\begin{pmatrix}0&1\\0&0\end{pmatrix},\qquad
 h=\begin{pmatrix}1&0\\0&-1\end{pmatrix},\qquad
 f=\begin{pmatrix}0&0\\1&0\end{pmatrix},
 \quad \mathfrak n=\F_5e,\quad\mathfrak b=\F_5e\oplus\F_5h.
\]
These are $G_L$-stable subspaces of $M=\mathfrak{sl}_2(\F_5)$.

\begin{lemma}[Restriction of the exceptional class]\label{lem:restriction}
Every nonzero class in $H^1(\SL_2(\F_5),M)$ has nonzero image in
$H^1(U,M/\mathfrak b)$, and hence in $H^1(L,M/\mathfrak b)$ when
$G_L\twoheadrightarrow U$ is the Tate residual representation.
\end{lemma}
\begin{proof}
The centre has order 2 and acts trivially on $M$, so we may equally work with
$A_5$. Restriction to its Sylow 5-subgroup is injective, because restriction
followed by corestriction is multiplication by $[A_5:U]=12$.
For $u=\left(\begin{smallmatrix}1&1\\0&1\end{smallmatrix}\right)$,
\[
 u(e)=e,\qquad u(h)=h-2e,\qquad u(f)=f+h-e.
\]
Thus $(u-1)M=\mathfrak b$ and $1+u+u^2+u^3+u^4=0$ on $M$.
The cyclic-group formula gives
\[
 H^1(U,M)=M/(u-1)M=M/\mathfrak b.
\]
Projection to the trivial module $M/\mathfrak b$ induces an isomorphism on
$H^1(U,-)$. Its inflation to $G_L$ is injective. The resulting character
is a nonzero scalar multiple of $a_q$.
\end{proof}

Let $\mathcal L_L\subset H^1(L,M)$ be the tangent space of ordinary,
fixed-determinant deformations relative to their weight algebra. Its local
Tate-dual annihilator is denoted by $\mathcal L_L^\perp$.

\begin{lemma}[The ordinary condition kills the exceptional line]
\label{lem:kill}
For the Tate representation satisfying \eqref{eq:tate},
\[
 \mathcal L_L=\im\bigl(H^1(L,\mathfrak n)\longrightarrow H^1(L,M)\bigr).
\]
Consequently a nonzero exceptional class inflated from
$H^1(\SL_2(\F_5),M)$ does not belong to $\mathcal L_L^\perp$.
\end{lemma}
\begin{proof}
In an upper-triangular ordinary deformation with fixed determinant and fixed
inertial weights, the diagonal infinitesimal character is unramified. Lifting
the extension requires its cup product with $a_q$ to vanish. The pairing of
$a_q$ with a nonzero unramified character is a nonzero multiple of $v_L(q)$.
Thus that diagonal direction is zero. All upper extension directions
$H^1(L,\mathfrak n)$ remain available. Changing the ordinary line only changes
the class by a coboundary. This proves the tangent-space assertion; it is also
the fixed-determinant form of AKT's calculation in Proposition 4.3.

For the trace pairing on $M$, one has $\mathfrak n^\perp=\mathfrak b$.
Functoriality and nondegeneracy of local Tate duality therefore imply
\[
 \mathcal L_L^\perp\subset
 \ker\bigl(H^1(L,M(1))\longrightarrow
                     H^1(L,(M/\mathfrak b)(1))\bigr).
\]
Here the mod-5 cyclotomic twist is trivial. Lemma~\ref{lem:restriction}
shows that the exceptional class has nonzero image under this map.
\end{proof}

\begin{remark}
The condition $5\nmid v_L(q)$ does two jobs. It ensures a ramified unipotent
residual image, and makes the ordinary deformation ring smooth relative to
weight. Nonzero unipotent image alone is not being substituted for both facts.
\end{remark}

\section{The ordinary local ring}
\label{sec:ring}

Let $d=[L:\Q_5]$, and let $\OO$ be a sufficiently large coefficient ring.
Put $\Lambda_L=\OO[[(\OO_L^\times)^f]]$, with $(\OO_L^\times)^f$ the maximal
torsion-free pro-5 quotient. Thus $\dim\Lambda_L=1+d$. We use the weight branch
through the Tate representation. The fixed determinant is $\e_5^{-1}$.

\begin{proposition}[Fixed-determinant form of AKT, Proposition 4.3]
\label{prop:smooth}
For \eqref{eq:tate}, the ordinary lifting ring is formally smooth over
$\Lambda_L$, of relative dimension $d+3$ and total dimension $2d+4$.
Moreover
\[
 \dim\mathcal L_L=d+1,\qquad h^0(L,M)=1,
 \qquad \dim\mathcal L_L-h^0(L,M)=d.
\]
\end{proposition}
\begin{proof}
AKT, Section 4.2.1, is stated for a general prime $p$ under exactly the
Kummer-valuation hypothesis \eqref{eq:tate}. Proposition 4.1 proves
representability of the ordinary condition and uniqueness of its ordered
ordinary line; Proposition 4.2 identifies it with the matrix relations provided
by ordinary local--global compatibility. Proposition 4.3 proves formal
smoothness, with two weight characters and without fixing the determinant.
The ring there has dimension $5+3d$ over a weight ring of dimension $1+2d$.

At the odd prime 5, first impose the determinant relation on the two inertial
weight characters. This removes $d$ weight variables. The remaining unramified
determinant character splits off by its unique square root congruent to 1.
Fixing it removes one further variable. This gives dimension $2d+4$ over
$\Lambda_L$, and relative dimension $d+3$.

For clarity, the essential infinitesimal argument can be given directly.
Write an ordinary lift as
\[
 \begin{pmatrix}\phi&*\\0&\e_5^{-1}\phi^{-1}\end{pmatrix},
\]
so its extension module has character $\e_5\phi^2$. For a small extension
of coefficient algebras with kernel $I$, the obstruction to lifting the
extension cocycle lies in $H^2(L,\F_5)\otimes I$, a one-dimensional
obstruction space. Changing the unramified part of $\phi$ by $x\ur$
changes that obstruction by
\[
 \pm2x(\ur\smile a_q).
\]
Its invariant is a nonzero multiple of $2xv_L(q)$, so every obstruction can
be removed. The ordinary line is recovered from an inertia matrix with unit
upper-right entry. The forgetful map from the flagged ordinary functor is
therefore an isomorphism at this residual representation.

At fixed weight the diagonal infinitesimal direction is forced to vanish.
The upper cocycles contribute $h^1(L,\F_5)=d+2$ parameters and the line
contributes one. The coboundary dimension is $3-h^0(L,M)=2$, giving
$\dim\mathcal L_L=d+1$. In particular the Kummer torsion directions have
not been discarded by taking a flat quotient.
\end{proof}

\section{The exceptional-five lifting argument}
\label{sec:patch}

\begin{proposition}[Proposed exceptional-5 ordinary lifting variant]
\label{prop:global}
Let $K$ be a CM field containing $\mu_5$. Let
$\rho:G_K\to\GL_2(\overline{\Q}_5)$ have determinant $\e_5^{-1}$,
be unramified almost everywhere, and be ordinary of weight zero. Suppose:
\begin{enumerate}
\item $\bar\rho$ is decomposed generic and has image $\SL_2(\F_5)$;
\item $\bar\rho$ is the residual representation of a cuspidal, ordinary,
weight-zero representation $\pi$ of $\PGL_2(\mathbb A_K)$;
\item at every place above 5, $\rho$ and $r_5(\pi)$ are the cohomology
representations of split Tate elliptic curves, with the common residual
Kummer class having valuation nonzero modulo 5. In particular their local
representations are extensions
$0\to\Q_5\to V\to\Q_5(-1)\to0$, after extending coefficients.
\end{enumerate}
The proposed lifting assertion $\Tinput$ is that $\rho$ is automorphic.
We give the patching adaptation below as a research draft of this assertion.
\end{proposition}

We give the full construction as a candidate proof, using the published
abstract patching theorem. The assertion is not an application of the excluded
case of AKT, Theorem A.14, or CN, Theorem 5.2.

\subsection{Preparation and the two deformation problems}

Make the soluble CM base changes used before AKT, Theorem A.7, but require
the extensions at every 5-adic place to be unramified. This preserves
\eqref{eq:tate}, and allows $[K_v:\Q_5]>4$. The imaginary quadratic auxiliary
fields can be chosen to split at 5 and at a protected decomposed-generic
rational prime. The latter prime then remains a witness after every extension.
Arrange that the automorphic seed is Iwahori-spherical and that, away from 5,
the Galois representations are unramified or unipotently ramified, with
trivial residual representation at the level-raising places. These are
local soluble extension requirements and can be imposed together with the
unramified requirements at 5. Full $\SL_2(\F_5)$ image survives soluble
restriction. The printed step which trivializes the residual representation
at 5 is omitted; Proposition~\ref{prop:smooth} supplies its replacement.

Write
\[
D=[K^+:\Q],\quad d_v=[K_v:\Q_5],\quad
\Lambda=\widehat\bigotimes_{v\mid5,\OO}\Lambda_v,
\qquad \dim\Lambda=1+2D.
\]
We use one inertial weight character at each $v$, with ordered diagonal
characters $\phi_v$ and $\e_5^{-1}\phi_v^{-1}$. The maximal torsion-free
pro-5 weight branch contains the Tate point: under reciprocity the cyclotomic
character is the local norm, up to inversion, and is trivial on the 5-power
roots of unity of $K_v$. Let
$\mathfrak p_0=\ker(\Lambda\to\OO)$ be the weight-zero augmentation.

Let $S=S_5\cup R$, with $R$ the prepared prime-to-5 ramification set. Use
the smooth fixed-determinant ordinary ring of Proposition~\ref{prop:smooth}
at every place above 5. At $v\in R$, use the fixed-determinant rings of AKT,
Lemma A.2, with tame characters $x_v=1$ or $x_v=\chi_v$, where $\chi_v$ has
5-power order and $\chi_v^2\ne1$. Such characters exist after the preparation
$q_v\equiv1\pmod5$ and enlarging the coefficient field. The resulting problems
are denoted $\mathcal S_1$ and $\mathcal S_\chi$; they are canonically identical
modulo the coefficient uniformizer $\varpi$. Additional smooth prime-to-5
places allowed by Lemma A.3 cause no change in the argument.

\subsection{The retained local condition and Taylor--Wiles primes}

Choose $v_0\mid5$ and put
\[
T=S\setminus\{v_0\},\qquad t=|T|,\qquad d_0=d_{v_0}.
\]
The split imaginary quadratic auxiliary extension gives at least two places
above 5, so $T$ is nonempty. In the $T$-framed relative dual Selmer group,
places of $T$ have no dual local condition, whereas $v_0$ retains
$\mathcal L_{v_0}^{\perp}$. The relative tangent is taken over the fixed
weight algebra; Proposition~\ref{prop:smooth} gives
\[
\dim\mathcal L_{v_0}=d_0+1,\qquad h^0(K_{v_0},M)=1.
\]
Lemma~\ref{lem:kill} excludes the inflated exceptional line from this dual
Selmer group. The local condition is the same for $\mathcal S_1$ and
$\mathcal S_\chi$. Taking $T=S$ would drop this condition and invalidate
the exclusion.

Here is the replacement for AKT, Lemma A.5. Put $K_N=K(\mu_{5^N})$.
The residual field is disjoint from the abelian extension $K_N/K$ because
$\SL_2(\F_5)$ is perfect. The combined finite image is
$\SL_2(\F_5)\times\Gal(K_N/K)$. Since the adjoint module has no
$\SL_2(\F_5)$-invariants, the $H^1$ of this combined image is precisely the
exceptional line. A dual Selmer class vanishing on the kernel of the combined
representation would inflate from that line, and hence must be zero.

Any remaining nonzero class restricts to a nonzero equivariant homomorphism
from this kernel to $M$. The span of its image over the residue coefficient
field is all of $M$, by absolute irreducibility of the adjoint module. This
span formulation also covers enlarged residue coefficient fields. Choose a
lift of $g=\diag(2,3)$ with trivial $5^N$-cyclotomic character. Varying it by
kernel elements makes the class nonzero in $M/(g-1)M$. Chebotarev supplies a
regular semisimple Taylor--Wiles prime of norm $1$ modulo $5^N$ detecting the
class. Degree-one primes may be chosen as in AKT, Lemma A.5; they have the
required imaginary quadratic splitting property.

Iteration and padding produce sets $Q_N$ of one fixed size $q$, with the
relative dual Selmer group zero for both problems. Set $Q_0=\varnothing$.
AKT, Proposition A.4, gives presentations with
\begin{equation}\label{eq:generators}
g=q-3D-1+t+d_0
\end{equation}
global variables over the partially framed local ring. Here the omitted
ordinary place contributes $d_0$, each Taylor--Wiles place contributes one,
and the complex places contribute $-3D$. Increase $q$ if necessary. Only the
presentation-generator bound is used: dual Selmer vanishing does not assert
formal smoothness of the augmented global ring.

\subsection{Framing algebras and finite-level presentations}

Write
\[
R_N=R_{\mathcal S_1,Q_N},\quad R'_N=R_{\mathcal S_\chi,Q_N},\quad
\Delta_N=\prod_{v\in Q_N}k(v)^\times(5).
\]
Let $\mathcal B=\Lambda[[z_1,\ldots,z_{4t-1}]]$ be the formal framing
algebra. Absolute irreducibility and $T\ne\varnothing$ identify the framed
rings, after compatible choices of representatives and coordinates, with
\[
R_{\mathcal S_1,Q_N}^{T}\simeq
\mathcal B\widehat\otimes_\Lambda R_N,
\qquad
R_{\mathcal S_\chi,Q_N}^{T}\simeq
\mathcal B\widehat\otimes_\Lambda R'_N.
\]
These choices can respect both the Taylor--Wiles augmentation and reduction
modulo $\varpi$, as in AKT, Section 6.2. Fixing the determinant of the
representation does not impose determinant-one framing matrices: the framing
count remains $4t-1$.

Let $\Lambda_T=\widehat\bigotimes_{v\in T,\OO}\Lambda_v$, taking
$\Lambda_v=\OO$ away from 5. The local coefficient rings are
\[
R^{\mathrm{loc}}=
\Lambda\widehat\otimes_{\Lambda_T}
 \left(\widehat\bigotimes_{v\in T,\OO}R_{v,1}\right),\qquad
R^{\prime\mathrm{loc}}=
\Lambda\widehat\otimes_{\Lambda_T}
 \left(\widehat\bigotimes_{v\in T,\OO}R_{v,\chi}\right).
\]
Thus the weight factor $\Lambda_{v_0}$ remains present, although its ordinary
lifting-ring factor is omitted. Restriction gives maps from these rings to
the corresponding $T$-framed global rings. Define
\begin{gather*}
R_\infty=R^{\mathrm{loc}}[[X_1,\ldots,X_g]],\qquad
R'_\infty=R^{\prime\mathrm{loc}}[[X_1,\ldots,X_g]],\\
\Delta_\infty=\Z_5^q,\qquad
S_\infty=\mathcal B[[\Delta_\infty]],\qquad
\mathfrak a_\infty=\ker(S_\infty\to\Lambda).
\end{gather*}

\subsection{The five finite-level patching inputs}

We verify the inputs of \cite{ACC}, Section 6.4.1. Its formal framing algebra
is arbitrary; it does not require framing at every ramified place.

\begin{enumerate}[label=\textbf{\arabic*.},leftmargin=24pt]
\item \textbf{Diamond groups.}
The group $\Delta_N$ is a quotient of $\Delta_\infty$ whose kernel is
contained in $5^N\Delta_\infty$, because every Taylor--Wiles norm is $1$
modulo $5^N$. At $N=0$ the quotient is trivial.

\item \textbf{Perfect complexes and their reductions.}
Take the complexes constructed before AKT, Proposition A.13, and specialize
the full diamond-operator weight algebra $\Lambda_1$ to $\Lambda$:
\[
C_N=F_{\mathfrak m,1,Q_N}\otimes_{\Lambda_1[\Delta_N]}\Lambda[\Delta_N],
\qquad
C'_N=F_{\mathfrak m,\chi,Q_N}\otimes_{\Lambda_1[\Delta_N]}\Lambda[\Delta_N].
\]
Here $\mathfrak m$ is the non-Eisenstein maximal ideal attached to the
common residual representation. AKT, Lemma A.10 and its Taylor--Wiles-level
construction give perfectness; the non-neat-level input is supplied by
Theorem 5.11 using $\mu_5\subset K$. The tame coefficient systems $1$ and
$\chi$ agree modulo $\varpi$, giving the required derived identification of
$C_N/\varpi$ and $C'_N/\varpi$. Their finite Hecke algebras, denoted
$\mathbb T_N$ and $\mathbb T'_N$, have the same image after this reduction.
Changing the framing set changes none of these complexes or operators.

\item \textbf{Global rings and local structures.}
The rings $R_N,R'_N$ have their usual $\Lambda[\Delta_N]$-structures and
agree modulo $\varpi$. The restriction maps just constructed give the
required $R^{\mathrm{loc}}$- and $R^{\prime\mathrm{loc}}$-structures on
$\mathcal B\widehat\otimes_\Lambda R_N$ and
$\mathcal B\widehat\otimes_\Lambda R'_N$. They respect the common reduction.

\item \textbf{Uniform presentations.}
Equation~\eqref{eq:generators} supplies surjections
\[
R_\infty\twoheadrightarrow\mathcal B\widehat\otimes_\Lambda R_N,
\qquad
R'_\infty\twoheadrightarrow\mathcal B\widehat\otimes_\Lambda R'_N.
\]
Choose generators compatibly with $N=0$ and modulo $\varpi$: fix generators
at level zero, lift them through the augmentation maps, and extend them by
generators in the augmentation kernels. Choose a common generating tuple
on the reduced rings and lift it to both sides. Nakayama's lemma and the
uniform bound give the required presentations.

\item \textbf{Deformation maps to the Hecke actions.}
The adaptations of AKT, Propositions A.11 and A.13 give surjections
\[
R_N\twoheadrightarrow\mathbb T_N/I_N,\qquad
R'_N\twoheadrightarrow\mathbb T'_N/I'_N,
\qquad I_N^\delta=(I'_N)^\delta=0
\]
for a single $\delta$ independent of $N$. To verify the changed local
condition, first adjoin the ordinary $U$-operators to the Hecke algebra.
AKT, Corollary 5.9 supplies the ordered characters, determinant
$\e_5^{-1}$, and the characteristic-polynomial and ordered-product matrix
identities. Proposition 4.2 identifies those identities with the actual
ramified ordinary lifting functor. Descent of the Galois representation to
the algebra generated by traces and weights, together with the descent
remark following Proposition 4.2, descends this representable local condition
to the smaller Hecke algebra modulo its intersected nilpotent ideal.
The determinant already has the required value. This is the argument of
AKT, Theorem 6.3 and Corollary 6.4, with its $\PGL_2$ form in Corollary 5.9.
The nilpotence bound stays uniform. Frobenius traces and the
weight/diamond structures determine the maps, hence give their compatibility
modulo $\varpi$. A genuine action of $R_N$ on the complex before passage
through nilpotent ideals is not being asserted.
\end{enumerate}

AKT, Lemma A.12 and the construction after it also give
\[
C_N\otimes^{\mathbf L}_{\Lambda[\Delta_N]}\Lambda\simeq C_0,
\qquad
R_N\otimes_{\Lambda[\Delta_N]}\Lambda\simeq R_0,
\]
with corresponding primed identities. Killing the diamond characters makes
the Taylor--Wiles local representations unramified. The identities respect
reduction modulo $\varpi$, the Hecke maps, and the chosen presentations;
these are the augmentation requirements of \cite{ACC}, Section 6.4.1.

\subsection{Patched objects and the support hypotheses}

The abstract construction now produces perfect complexes $C_\infty,C'_\infty$
over $S_\infty$, finite Hecke algebras in their derived endomorphism rings,
and compatible surjections from $R_\infty,R'_\infty$ modulo uniformly
nilpotent ideals. It supplies the $S_\infty$-algebra structures on the patched
rings. ACC+, Proposition 6.4.17 gives the identified reductions and their
compatible ring actions. Lemma 6.4.15 and Proposition 6.4.16 give
\begin{equation}\label{eq:augmentation}
R_\infty/\mathfrak a_\infty\twoheadrightarrow R_0,
\qquad
C_\infty\otimes^{\mathbf L}_{S_\infty}S_\infty/\mathfrak a_\infty
\simeq C_0,
\end{equation}
and compatibility with the Hecke maps. Only the indicated surjection of
rings is needed.

Here are all three hypotheses of ACC+, Assumption 6.3.6.

\paragraph{Dimensions.}
Since $\sum_{v\mid5}d_v=2D$ and each prime-to-5 local factor has relative
dimension three, one has
\begin{align*}
\dim R^{\mathrm{loc}}&=1+4D-d_0+3t,\\
\dim R_\infty&=D+4t+q,\\
\dim S_\infty&=2D+4t+q.
\end{align*}
The same formulas hold for the primed ring. The smooth factors are
$\OO$-flat. The dimension and component statements of AKT, Lemma A.2,
together with the completed tensor-product calculation, give
\[
\dim(R_\infty/\varpi)=\dim R_\infty-1,\qquad
\dim(R'_\infty/\varpi)=\dim R'_\infty-1.
\]
Hence
\begin{equation}\label{eq:defect}
\dim S_\infty-\dim R_\infty=D.
\end{equation}
Relative to the printed $T=S$ calculation, deleting $R_{v_0}$ removes
$d_0+3$ variables over weight, while \eqref{eq:generators} adds $d_0-1$.
Both patched dimensions therefore fall by four.

\paragraph{Components.}
The retained 5-adic local factors are formal power series rings over their
weight factors; the omitted place contributes its smooth weight algebra.
They introduce no extra components. The prime-to-5 level-raising factors
are exactly those in AKT, Lemma A.2. Its first part gives the unique generic
component above each maximal generic special-fibre component on the
unprimed side; its second part gives the unique characteristic-zero generic
component on the primed side. The completed tensor-product and power-series
argument at the end of AKT, Theorem A.7 applies with these same factors.
The remaining factors are smooth, so the required component and
specialization conditions hold. In this all-Tate problem the factors are
equidimensional; every characteristic-zero target point lies on a component
of maximal dimension.

\paragraph{Generic cohomological concentration.}
The prime $\mathfrak p=(\mathfrak a_\infty,\mathfrak p_0)$ of $S_\infty$
has one-dimensional characteristic-zero quotient. By \eqref{eq:augmentation},
specialization there and inversion of 5 gives the weight-zero ordinary
localized $\PGL_2$ complex at level zero. The seed $\pi$ supplies nonzero
rational cohomology. AKT, Theorem 5.10 concentrates this cohomology in an
interval of length $D$, up to the degree reversal for the chain/dual convention.
This is precisely Assumption 6.3.6(3) with $l_0=D$; no integral amplitude
vanishing is assumed.

\subsection{The target point and automorphy}

The target representation defines $f:R_0\to\OO$ at weight zero. Pull
$\ker f$ back under the surjection in \eqref{eq:augmentation}, and denote
the resulting prime of $R_\infty$ by $\mathfrak x$. Its contraction to
$S_\infty$ is $(\mathfrak a_\infty,\mathfrak p_0)$ and it lies on a maximal
component. ACC+, Proposition 6.3.8 and Corollary 6.3.9 therefore place it in
the support of
\[
H^*\!\left(C_\infty\otimes^{\mathbf L}_{S_\infty}
S_\infty/(\mathfrak a_\infty,\mathfrak p_0)\right)[1/5].
\]
Proposition 6.4.16 and the compatible Hecke maps transfer this support to
the rational weight-zero ordinary complex at level $U(1,1)$. By AKT,
Theorem 5.10(3), the corresponding eigensystem comes from a cuspidal ordinary
weight-zero automorphic representation. Frobenius traces and absolute
irreducibility identify its Galois representation with $\rho$.
Finally \cite{ACC}, Proposition 6.5.13 descends automorphy through the
prepared soluble extension, whose 5-adic ramification was explicitly
controlled. This completes the candidate proof of $\Tinput$.

\section{A modular seed preserving the Tate condition}
\label{sec:seed}

The following is a local-preservation variant of AKT, Proposition 9.13.
It uses AKT's unconditional seed construction, independently of
Proposition~\ref{prop:global}.

\begin{lemma}[Preserving the completions at 5]\label{lem:seed}
Let $A/K$ be an elliptic curve over a CM field, with split multiplicative
reduction at every $v\mid5$ and local Tate parameters $q_v$.
Put $\bar r=A[5]$, with its Weil pairing. Then, after a solvable CM extension which
splits completely at the 5-adic places, there is a modular elliptic curve
$B$ with a symplectic isomorphism $B[5]\simeq\bar r$.
Its Tate parameters at 5 can be chosen in
prescribed sufficiently small neighbourhoods of the $q_v$. A finite
avoidance field and a chosen decomposed-generic prime can also be retained.
\end{lemma}
\begin{proof}
In AKT's proof of Proposition 9.13, the initial solvable CM extension
trivializes the prescribed residual representation at 2, 3 and 5. Triviality
at 5 is used to supply local Tate points on the twist $Y_{\bar r}$ of
$Y(5)$. Here the local curves $A_{K_v}$ with their tautological level
structures already provide points on this fixed symplectic twist.
Trivialize only at 2 and 3,
arrange their splitting over the maximal real subfield and the required
unramified $i$-extension at 2, and keep the places above 5 completely split.

The curve $Y_{\bar r}$ is a nonempty open of a projective line. AKT,
Lemma 9.7, therefore applies with the prescribed local 5-adic opens, as
well as the opens at 2 and 3 used in Proposition 9.13. It gives an elliptic
curve $B$ with full and decomposed-generic mod-3 image, the required
discriminant conditions at 2 and 3, and the specified mod-5 representation.
AKT, Proposition 9.12, proves $B$ modular. That proposition imposes no
condition at 5 and no exclusion of $\zeta_5$.

The internal extension used to prove $B$ modular need not preserve the
5-adic parameters: modularity is descended back to the field on which
$B$ was constructed before our proposed 5-adic lifting argument is used.
\end{proof}

The preservation condition is on the local moduli points of
$Y_{\bar r}$, including their symplectic level structures. In particular
the approximation takes place on the prescribed residual twist.

\section{Tate specialization of the geometric family}\label{sec:tategeometry}
\label{sec:geometry}

Return to the covers of family 030. After pulling the parameter line back
to the fixed elliptic curve $E$, write the moving point as $R$ and take
$n=3p$. Choose a point $R_*$ of exact order 6. Then
\[
 [n]R_*=[3p]R_*\in E[2]\setminus\{0\},\qquad
 -2(n/3)R_*=-2pR_*
\]
has exact order 3. This is enough for the degree-3 version of the second
degeneration: full $6p$-torsion is unnecessary.

For the normalized covering function on the constant-$E$ parameter cover,
write $f_{j,R}$ for the rational function, over the field obtained by adjoining
$R$, with
\[
\operatorname{div}(f_{j,R})=([j]R)-(-[j]R)-j(R)+j(-R),
\qquad f_{j,R}(O)=1.
\]
These are the functions of \srcM, Lemma 3.1. They satisfy the exact identity
\[
 f_{ab,R}=f_{a,[b]R}\,f_{b,R}^{\,a}.
\]
It follows by comparing divisors and the normalization $f(O)=1$.
For $a=3$ and $b=p$, the change $Y\mapsto Y/f_{p,R}$ identifies
the degree-3 cover for $n=3p$ with the degree-3 family for $n=3$
pulled back by $[p]$. The identity is compatible with $w$.
Since $[p]R_*=\pm R_*$, the boundary splitting data can be fixed
independently of $p$.

\begin{lemma}[Width three at the elliptic endpoint]\label{lem:width}
After a fixed extension of a characteristic-zero local ground field
splitting the degeneration data, the genus-one quotient
$C_3/\langle w\rangle$ has an $I_3$ degeneration at $R_*$. For a local
parameter $s$ there,
\begin{equation}\label{eq:qexp}
 q_{A}(s)=a s^3(1+O(s)),\qquad a\ne0.
\end{equation}
In particular, at a place of residue characteristic 5, arbitrarily small
nonzero specializations can satisfy $5\nmid v(q_A(s))$.
\end{lemma}
\begin{proof}
The degree-3 special cover consists of a connected unramified elliptic cover
$E'\to E$ and a rational component joined by three nodes, each with
smoothing order one. The connectedness follows from the exact-order-3
divisor class $-2pR_*$. The involution $w$ preserves both components,
acts by inversion on the elliptic component, and acts on the three node
labels by $i\mapsto-i$. There is one fixed node and one pair.

At the fixed node, $w$ acts by $(-1,-1)$ on the two branch parameters.
If $xy=s$ is the smoothing equation, the quotient has equation
$XY=s^2$. The pair of nodes yields one quotient node of thickness one.
The two quotient components are rational. Resolving the thickness-two
node inserts one component, yielding the three-component cycle $I_3$.
Thus the quotient Jacobian has $j$-invariant with a pole of order 3,
which gives \eqref{eq:qexp}.

For the splitting field, it is enough to adjoin $E[6]$. The point
$(O,1)$ makes $E'\to E$ a 3-isogeny whose kernel $\mu_3$ is rational.
Its restriction $E'[2]\to E[2]$ is an isomorphism of Galois modules.
If $Q=[n]R_*$, its unique 2-torsion lift $Q'$ is therefore rational,
and the three attaching points are $Q'+\ker(E'\to E)$.
They are all rational. The rational tail and the quotient cycle are
then split as well.

In the arithmetic local field, the constant $a$ need not be an integer
unit. For sufficiently small $s$, however,
\[
 v(q_A(s))=v(a)+3v(s).
\]
Choose arbitrarily large integers $k=v(s)$ avoiding the single residue
class $-v(a)/3$ modulo 5. Splitness is preserved in a sufficiently small
neighbourhood, for example by the square condition on $-c_6$ in a local
Weierstrass model. Finally $t=x(R)$ is etale at the exact-order-6 point
$R_*$, so these conditions define nonempty opens in the original
parameter line over the extended local field.
\end{proof}

The extension splitting the local boundary data must be chosen before the
parameter is specialized. Taking an arbitrary extension afterwards could
multiply the valuation by 5 and invalidate \eqref{eq:tate}. Such local
extensions can be accommodated in a solvable CM extension. Here one can
make the approximation explicit: take, at each 5-adic place of $K^+$, the
Galois compositum containing the relevant completions of $K$, their
6-torsion and their conjugates. The preceding proof shows this already
splits the boundary. Its Galois group has order divisible only by 2 and 3,
so it is solvable. Realize successive cyclic steps by the degree-2 and
degree-3 cases of Grunwald--Wang, keeping every real place and any
prescribed auxiliary places split. Taking a Galois closure preserves
total reality, solvability, and the prescribed complete splitting.
Composing that totally real extension with $K$ gives the required CM
extension. The local character theorem used here is recorded, for
example, in Conrad, Proposition A.3~\cite{Conrad}; there is no Wang
exception for the cyclic steps of degree 2 or 3.

There is no circularity in the choice of $p$: the 6-torsion preparation
is independent of $p$, so one can fix the solvable CM field first and
then choose $p$ large enough for the original residual-image estimates.
The variable character pieces use the
independent-factor Hilbert argument over the fixed field, not a new
degree-dependent union bound for $p$.

\section{Completion of the CM argument}
\label{sec:consequence}

The geometric package is $\Ginput$ from Section~\ref{sec:package};
the additional split Tate specialization is Lemma~\ref{lem:width}.
We now combine it with $\Tinput$ from Proposition~\ref{prop:global}.

\begin{corollary}[Consequence of the proposed lifting variant]
\label{cor:allcm}
If Proposition~\ref{prop:global} is validated, then
$\Ginput$ implies the modularity of every elliptic curve over
every CM number field, including fields containing $\zeta_5$.
\end{corollary}
\begin{proof}
The case $\zeta_5\notin K$ is Theorem~\ref{thm:CM}. For the remaining
case, elliptic curves with complex multiplication are treated by Hecke
characters. For a non-CM curve, first take a solvable CM extension admitting
the 6-torsion local boundary data and then choose a large $p>5$ with the
residual properties needed for \eqref{eq:switch}. Impose:
\begin{itemize}
\item the compact-type local opens at places above $3p$, as in the original
switch, so that its successive representations have matching local types;
\item at every place above 5, the second-degeneration opens of
Lemma~\ref{lem:width}, so that $A_t$ is Tate and
$5\nmid v(q_{A_t})$;
\item the finite residual-image and decomposed-generic conditions on
$A_t[p]$, $A_t[5]$, and the required order-$p$ character pieces modulo 3.
\end{itemize}
Weak approximation together with Hilbert irreducibility gives a common
specialization. The 5-adic conditions do not conflict with the opens at
3 and $p$, or with full global mod-5 monodromy.

Apply Lemma~\ref{lem:seed} to $A_t[5]$, preserving the 5-adic completions,
the residual avoidance fields and decomposed-generic primes. It supplies a
modular $B$ with the same residual representation and Tate type at 5.
Proposition~\ref{prop:global} would then prove $A_t$ modular. Solvable
descent returns this modularity to the field used for specialization.

Now pass to $F=K'(\mu_{3p})$, where $K'$ denotes that field. The three
ordinary/potentially Barsotti--Tate lifting steps in \eqref{eq:switch}
use CN, Theorem 5.2, only at $p,3,p$. Its exceptional-5 condition never
enters these steps. They give the modularity of $B_t$, hence of $E$ over
$F$. Abelian cyclotomic descent and then solvable descent give modularity
over the original $K$.
\end{proof}

\section{An independent residual-mod-three consequence}

The proposed lifting variant has a consequence that does not use
$\Ginput$.

\begin{corollary}[Conditional extension of CN, Theorem 6.1(1)]\label{cor:mod3}
Assume Proposition~\ref{prop:global}. Let $F$ be any CM number field and
let $E/F$ be an elliptic curve. If $E[3]$ is decomposed generic and
$E[3]|_{G_{F(\zeta_3)}}$ is absolutely irreducible, then $E$ is modular.
In particular this conclusion would no longer require $\zeta_5\notin F$.
\end{corollary}
\begin{proof}
Only $\mu_5\subset F$ is new. In the proof of AKT, Proposition 9.15,
and CN, Proposition 6.1.6, first make the soluble CM preparation that
trivializes the prescribed mod-3 representation at the places above 5.
The local fields then contain $\mu_3$, by the determinant condition.
On the prescribed twist of $Y(3)$, choose at such a place a Tate curve
with parameter
\[
 q=a^3,\qquad v(a)>0,\qquad 5\nmid v(a).
\]
Its mod-3 representation is trivial, and $5\nmid v(q)$. These conditions
hold on a nonempty local open of the twist. The prescribed ordinary,
supersingular, or Tate opens at 3 in CN, Proposition 6.1.6, can be
retained independently. AKT, Lemma 9.7, constructs a curve $C$ satisfying
all these opens, with full and decomposed-generic mod-5 image, including
the required finite avoidance condition.

Apply Lemma~\ref{lem:seed} to $C[5]$ and then
Proposition~\ref{prop:global} over the resulting soluble CM extension,
preserving its 5-adic completions, full mod-5 image and generic prime.
Descend the modularity of $C$ to the field on which it was first
constructed. This removes the root-of-unity restriction from the mod-3
seed construction without changing its avoidance or prescribed 3-adic
local-type conclusions. Finally CN, Theorem 5.2, at $p=3$ proves the
modularity of the original $E$, exactly as in Theorem 6.1(1).
The irreducibility and decomposed-generic hypotheses of that theorem
are retained; this is not yet all-CM modularity by itself.
\end{proof}


\section{Scope of the new argument}

The field extension in Sections~\ref{sec:package}--\ref{sec:TR} uses the source package $\Ginput$ and
published automorphy theorems. The extra arithmetic use of the second
geometric degeneration is its width three, which is invertible modulo five:
\begin{align*}
\text{width-three degeneration at five}
&\Longrightarrow \text{Tate parameter of valuation prime to five}\\
&\Longrightarrow \text{smooth ordinary local condition}\\
&\Longrightarrow \text{exclusion of the exceptional adjoint class}.
\end{align*}
The condition is imposed on the auxiliary elliptic curve. The original curve
need not be ordinary or multiplicative at five. The local ring is an existing
input; retaining its dual local condition in the partially framed patching
problem is the new global argument.

Independent review should concentrate on Proposition~\ref{prop:global},
especially the partially framed patched complexes, their Galois action through
the Hecke algebra, and the support comparison. The local cohomology, local-ring
calculation, preservation of the Tate class through the modular seed, and
width-three specialization have explicit proofs above. A separate verification
of $\Ginput$ remains necessary to certify the source manuscript itself.
These dependencies are retained in each main statement.

\begin{thebibliography}{99}
\bibitem[M]{M}
\texttt{openai/math}, family 030,
\emph{Modularity of elliptic curves over imaginary quadratic fields},
manuscript dated 4 October 2026.
\href{https://github.com/openai/math/blob/main/preprints/Modularity-of-elliptic-curves-over-imaginary-quadratic-fields-October-4-2026/paper.pdf}{GitHub manuscript PDF}.
Source locations: Sections 3--5, 7--9.

\bibitem[CN]{CN}
A. Caraiani and J. Newton,
\emph{On the modularity of elliptic curves over imaginary quadratic fields},
arXiv:2301.10509v3, 27 March 2025.
\url{https://arxiv.org/html/2301.10509v3}.
Theorem 5.2, Theorem 6.1, and Section 6.1.

\bibitem[AKT]{AKT}
P. B. Allen, C. Khare, and J. A. Thorne,
\emph{Modularity of $\GL_2(\F_p)$-representations over CM fields},
Cambridge Journal of Mathematics \textbf{11} (2023), 1--158.
Version used: arXiv:1910.12986v2, 2 September 2022.
\url{https://arxiv.org/abs/1910.12986}.
Sections 4.2.1, 5.2, 8--9, and Appendix A.

\bibitem[ACC]{ACC}
P. B. Allen, F. Calegari, A. Caraiani, T. Gee, D. Helm,
B. V. Le Hung, J. Newton, P. Scholze, R. Taylor, and J. A. Thorne,
\emph{Potential automorphy over CM fields},
Annals of Mathematics \textbf{197} (2023), 897--1113.
\url{https://arxiv.org/abs/1812.09999}.
Sections 6.3--6.4 and Proposition 6.5.13.

\bibitem[Serre]{Serre}
J.-P. Serre, \emph{Propri\'et\'es galoisiennes des points d'ordre fini des
courbes elliptiques}, Inventiones Mathematicae \textbf{15} (1972), 259--331.
\href{https://www.college-de-france.fr/media/jean-pierre-serre/UPL5874918517843398173_Serre_proprie_te_s_galoisiennes_des_courbes_elliptiques.pdf}{Author's collected-paper copy}.

\bibitem[AC]{AC}
J. Arthur and L. Clozel, \emph{Simple Algebras, Base Change, and the Advanced
Theory of the Trace Formula}, Annals of Mathematics Studies 120,
Princeton University Press, 1989.
\url{https://www.math.toronto.edu/arthur/pdf/30.pdf}.
Chapter 3, Theorems 4.2 and 5.1, used through \cite{ACC}, Proposition 6.5.13.

\bibitem[Conrad]{Conrad}
B. Conrad, \emph{Lifting global representations with local properties},
Appendix A, Proposition A.3 (Grunwald--Wang).
\url{https://math.stanford.edu/~conrad/papers/locchar.pdf}.
\end{thebibliography}
\end{document}
